\documentclass[11pt,a4paper]{article}

\usepackage[margin=2.4cm]{geometry}
\usepackage{fontspec}
\usepackage{booktabs}
\usepackage{array}
\usepackage{parskip}
\usepackage{titlesec}
\usepackage{xcolor}
\usepackage{enumitem}
\usepackage{microtype}
\usepackage[hidelinks]{hyperref}

\setmainfont{Latin Modern Roman}
\setsansfont{Latin Modern Sans}

\titleformat{\section}{\large\bfseries}{\thesection}{0.8em}{}
\titlespacing*{\section}{0pt}{1.4em}{0.5em}
\titleformat{\subsection}{\normalsize\bfseries}{\thesubsection}{0.7em}{}
\titlespacing*{\subsection}{0pt}{1.0em}{0.35em}

\setlist[itemize]{leftmargin=1.4em,itemsep=0.15em,topsep=0.3em}
\setlist[enumerate]{leftmargin=1.6em,itemsep=0.15em,topsep=0.3em}

\newcommand{\code}[1]{\texttt{\small #1}}

\begin{document}

\begin{center}
{\Large\bfseries Baseline Retrieval vs.\ Cross-Encoder Re-Ranking}\\[0.35em]
{\normalsize A comparison of two retrieval pipelines over the Magic: The Gathering Comprehensive Rules}\\[0.6em]
{\small LMA Assignment \quad$\cdot$\quad 20 October 2026}
\end{center}

\vspace{0.4em}

\section{Setup}

I used the \emph{Magic: The Gathering Comprehensive Rules} (effective 7 June 2024) as the
document. It is 292 pages of numbered rules, and it is full of the things this assignment is
looking for: prohibitions (``can't'', ``may not''), exceptions (``except'', ``unless''), and
the same words reused in unrelated contexts. The word \emph{counter}, for instance, means both
``counter a spell'' and ``a counter on a permanent'', which turned out to be a useful source of
keyword traps.

Rather than splitting the text every $N$ characters, I chunked on rule boundaries. The PDF is
parsed into individual numbered rules (\code{701.17a}, \code{704.5j}, and so on), the
table-of-contents pages are dropped, and consecutive rules are packed into chunks of about 700
characters with one rule of overlap. This gave 1{,}937 chunks, averaging 642 characters and
capped at 1{,}099. The cap matters because the cross-encoder truncates at 512 tokens, and a
chunk that gets silently truncated is a chunk the re-ranker cannot score properly.

Both pipelines share the same chunks, the same top-10 candidate list, and the same prompt
string. The prompt is defined once as a constant, so the two pipelines cannot drift apart. The
LLM was \code{google/gemini-2.5-flash} through OpenRouter at temperature 0, so any difference
between Answer~1 and Answer~2 comes from the retrieved context and not from sampling.

\section{Architecture comparison}

A bi-encoder puts the query and the chunk through the transformer \emph{separately} and gets
two fixed vectors back, which are then compared with cosine similarity. The important
consequence is that the chunk vector never depends on the query, so the whole corpus can be
embedded once, offline, and reused forever. Answering a query costs one embedding plus a vector
lookup.

A cross-encoder concatenates the query and the chunk into a single sequence and runs them
through the model \emph{together}. Every query token can attend to every chunk token, so the
model is judging the relationship between the two rather than comparing two summaries that were
produced in isolation. The output is one relevance score.

That is the whole trade-off. The cross-encoder sees more and scores better, but nothing can be
precomputed, because the score depends on the pair. Every candidate needs its own forward pass
at query time. The bi-encoder is cheap and cacheable but has to compress a chunk into a single
vector before it has any idea what the query will be, and that compression is where the failures
in the next section come from.

\section{Rank shifts and answer quality}

I did not pick the test queries by guessing. I wrote about a hundred candidate questions, found
the gold chunk for each by searching the corpus for the rule number that actually answers it,
and then recorded where that chunk landed in the baseline top-10 and where it landed after
re-ranking. The table below is the subset that made the final cut.

\begin{center}
\small
\begin{tabular}{@{}llcc@{}}
\toprule
\textbf{Category} & \textbf{Query (abridged)} & \textbf{Baseline} & \textbf{Re-ranked} \\
\midrule
Keyword trap   & Look at another player's hand?              & \#8 & \#1 \\
Keyword trap   & Target a creature with protection?          & \#9 & \#3 \\
Keyword trap   & Tap an already-tapped permanent?            & \#5 & \#1 \\
Exception      & Both win and lose simultaneously?           & \#7 & \#1 \\
Exception      & Mana ability while casting a spell?         & \#8 & \#2 \\
Exception      & Trample blocked by protection?              & \#5 & \#2 \\
Negation       & Respond to a mana ability?                  & \#8 & \#1 \\
Negation       & Sacrifice a permanent in the graveyard?     & \#4 & \#1 \\
Buried chunk   & Planeswalker at zero loyalty?               & \#8 & \#2 \\
Buried chunk   & Creature with lethal damage marked?         & \#5 & \#2 \\
\midrule
Control        & What happens when a spell is countered?     & \#1 & \#4 \\
\bottomrule
\end{tabular}
\end{center}

Ten of the eleven queries moved the gold chunk from rank 4--10 into the top three. The pattern
is consistent: whenever the query shares surface vocabulary with an irrelevant passage, the
bi-encoder's single pooled vector gets pulled toward the lexical match. ``Look at another
player's hand'' is about information access, but the top baseline chunk is about hand size,
because the word \emph{hand} is doing all the work.

The answer-level differences were more interesting than the rank numbers.

\textbf{The clearest failure was a misstatement.} For ``what happens if a player would both win
and lose the game at the same time'', the baseline put the correct rule (104.3f) at rank 7 and
instead surfaced the general draw rules. Answer~1 confidently said \emph{the game is a draw},
citing 104.4a. That is simply wrong: 104.3f says the player \emph{loses}. Answer~2, with 104.3f
in the top three, got it right. This is the case I would point to if someone asked why
re-ranking is worth the cost, because the baseline answer is not vague or hedged, it is
confidently incorrect.

\textbf{The second was an omission.} For ``can a player look at another player's hand?'',
Answer~1 never found 402.3. It reasoned from rule 811.5 (teammates cannot review each other's
hands) and rule 720.4 (a player controlling another player can see their hand) and produced a
hedged answer about ``specific conditions''. It is not exactly wrong, but it is not the rule
either. Answer~2 quotes 402.3 directly: a player cannot look at another player's hand but may
count the cards.

\textbf{Two queries produced outright refusals.} ``Can you respond to a mana ability?'' and the
trample-versus-protection question both returned ``the context does not contain the answer''
under the baseline, and both were answered correctly after re-ranking. The mana ability case is
a good illustration of the negation problem: the bi-encoder ranked chunks that \emph{mention}
mana abilities and priority above the chunk that states the prohibition (605.3b), so the model
never saw the sentence saying it cannot be responded to.

\textbf{One case got worse.} ``What happens when a spell is countered?'' starts with the gold
chunk at rank 1 and re-ranking pushes it to rank 4. The cross-encoder promoted a chunk about
abilities that trigger when a spell is countered, and another about resolution order, above the
chunk that actually defines countering. Answer~1 was correct; Answer~2 refused. I think the
explanation is that \code{ms-marco-MiniLM-L-6-v2} is trained on MS MARCO web passages, where a
passage that discusses a topic in context often looks more relevant than a terse definitional
one. It is optimised for ``is this passage relevant'', not ``is this the authoritative
definition''. I am keeping this case in the report rather than dropping it, because it shows the
method is not a free win.

\section{Latency}

Both models were warmed up with a dummy call first, then each query was timed 15 times. All
measurements are on CPU (12 cores, no GPU).

\begin{center}
\small
\begin{tabular}{@{}lrr@{}}
\toprule
\textbf{Stage} & \textbf{Mean (ms)} & \textbf{Std (ms)} \\
\midrule
Query embedding (bi-encoder)          & 30.79  & 9.18 \\
FAISS search (top-10)                 & 1.65   & 1.01 \\
\textbf{Baseline retrieval total}     & \textbf{32.44} & \textbf{9.94} \\
Cross-encoder scoring (10 pairs)      & 708.31 & 122.70 \\
\textbf{Re-ranking pipeline total}    & \textbf{740.75} & \textbf{119.28} \\
\bottomrule
\end{tabular}
\end{center}

Building the index took 131 seconds, but that is a one-off cost and is not part of the query
path. The cross-encoder accounts for about 96\% of the re-ranking pipeline and adds roughly
21.8$\times$ the baseline retrieval cost, at about 70.8\,ms per pair. The FAISS lookup itself is
negligible at 1.65\,ms, so the vector search is not the bottleneck at this corpus size.

In absolute terms the overhead is under a second, which is why re-ranking ten candidates is
perfectly practical. It is only impractical when you try to do it over the whole corpus, which
is the next question.

\section{Why not run the cross-encoder over everything?}

Because the score depends on the pair, there is nothing to precompute. Running the cross-encoder
over the full corpus means $N$ forward passes per query, where $N$ is the number of chunks.
Using the measured per-pair cost of 70.8\,ms:

\begin{center}
\small
\begin{tabular}{@{}lr@{}}
\toprule
\textbf{Corpus size} & \textbf{Cross-encoder cost per query} \\
\midrule
1{,}937 chunks (this corpus) & $\approx$ 137 seconds \\
100{,}000 chunks             & $\approx$ 2 hours \\
1{,}000{,}000 chunks         & $\approx$ 20 hours \\
\bottomrule
\end{tabular}
\end{center}

And none of it can be cached, because the query changes every time. The bi-encoder does not have
this problem: it embeds the corpus once (131 seconds here) and then answers each query in about
32\,ms.

This is the entire argument for retrieve-then-rerank. The cheap model cuts $N$ down to $k = 10$,
and the expensive model only ever sees those ten candidates. The re-ranking cost then becomes
independent of corpus size: it is about 708\,ms whether the corpus has two thousand chunks or
two million. The cross-encoder is not a replacement for the vector search, it is a second pass
that only makes sense after the first pass has done the filtering.

\section{Closing note}

The one thing I would change if I did this again is the prompt. My first version told the model
to say ``the provided context does not contain the answer'' when it could not find one, and that
instruction made it bail out whenever a distractor chunk was present, even when the gold chunk
was sitting right there in the context. Removing the escape hatch and asking it to cite the rule
number it used fixed most of the false refusals. It is a small reminder that in a RAG system the
retrieval and the prompt are not independent, and a prompt that is too cautious can hide a
retrieval improvement that is actually working.

\end{document}
